\notechapter{N-20220721}{Conditioning as Projection: Probability, Jensen's Inequality, and Lost Information}

Conditioning changes the weights assigned to outcomes. When the condition is a single event this is a new probability measure; when the available information is a partition of the sample space it is a new function, the conditional expectation. That function turns out to be an orthogonal projection in \(L^2\). The same operation also satisfies Jensen's inequality and, in information theory, can only decrease the relative entropy between distributions after outcomes have been merged.\footnote{\href{https://web.stanford.edu/class/stats310a/lnotes.pdf}{Stanford Stats 310A}, \S4.3 (conditional expectation as projection); MIT 6.441, Chapters 2--3 (information-theoretic data processing).}

\section{How Conditioning Reweights Outcomes}

On a finite probability space \(\Omega\), let \(p(\omega)>0\) denote the weight of each outcome. For an event \(B\subseteq\Omega\) with \(\Pr(B)>0\), the conditional law is
\[
 p_B(\omega)=\frac{p(\omega)\mathbf1_B(\omega)}{\Pr(B)}.
\]
Thus for a random variable \(X\),
\[
 \mathbb E[X\mid B]
 =\sum_{\omega\in\Omega}X(\omega)p_B(\omega)
 =\frac{\mathbb E[X\mathbf1_B]}{\Pr(B)}.
\]
Conditioning is an explicit reweighting of the same outcomes.

Bayes' theorem follows by comparing two ways of multiplying conditional and marginal probabilities:
\[
 \Pr(H\mid E)=\frac{\Pr(E\mid H)\Pr(H)}{\Pr(E)}.
\]
The odds form is especially informative:
\[
 \boxed{\frac{\Pr(H\mid E)}{\Pr(H^c\mid E)}
 =\frac{\Pr(E\mid H)}{\Pr(E\mid H^c)}
  \frac{\Pr(H)}{\Pr(H^c)}.}
\]
It says that evidence multiplies prior odds by a likelihood ratio. For a rare event with prior probability \(1/100\), a test of sensitivity \(9/10\) and false-positive rate \(1/10\) gives
\[
 \Pr(H\mid +)
 =\frac{(9/10)(1/100)}
 {(9/10)(1/100)+(1/10)(99/100)}
 =\frac1{12}.
\]
High sensitivity alone cannot overcome very small prior odds and a substantial false-positive rate.

\section{Observing a partition gives a random estimate}

Suppose the observation only tells us which cell of a partition
\(\Omega=C_1\sqcup\cdots\sqcup C_m\) contains the outcome. If \(\Pr(C_j)>0\), define
\[
 M(\omega)=\frac{\sum_{\eta\in C_j}p(\eta)X(\eta)}
                  {\Pr(C_j)}
 \qquad(\omega\in C_j).
\]
The function \(M\) is the conditional expectation of \(X\) given this observation. It is constant on each information cell and averages \(X\) using the conditional weights within that cell.

For example, take four outcomes of probabilities
\[
 (p_1,p_2,p_3,p_4)=\frac1{10}(1,2,3,4),
\]
with \(X=(0,2,3,7)\). If the observation distinguishes only \(\{1,2\}\) from \(\{3,4\}\), the best cellwise predictions are
\[
 M|_{\{1,2\}}=\frac{1\cdot0+2\cdot2}{1+2}=\frac43,
 \qquad
 M|_{\{3,4\}}=\frac{3\cdot3+4\cdot7}{3+4}=\frac{37}{7}.
\]
The mean is preserved:
\[
 \mathbb E[M]
 =\frac3{10}\frac43+\frac7{10}\frac{37}{7}
 =\frac{41}{10}=\mathbb E[X].
\]

\section{Conditional Expectation as Projection}

Equip real-valued functions on \(\Omega\) with the inner product
\[
 \langle X,Z\rangle=\mathbb E[XZ].
\]
The functions constant on each information cell form a vector subspace \(V\) of this finite-dimensional Hilbert space. For any \(Z\in V\), say \(Z=z_j\) on \(C_j\),
\[
\begin{aligned}
 \mathbb E[(X-M)Z]
 &=\sum_jz_j\sum_{\omega\in C_j}
       p(\omega)(X(\omega)-M(\omega))\\
 &=0.
\end{aligned}
\]
Thus the residual \(X-M\) is orthogonal to every function measurable from the observation:
\[
 \boxed{M=\operatorname{Proj}_V X.}
\]

For any other predictor \(Z\in V\), Pythagoras gives
\[
 \mathbb E[(X-Z)^2]
 =\mathbb E[(X-M)^2]+\mathbb E[(M-Z)^2].
\]
It follows that \(M\) uniquely minimizes the mean squared prediction error among all predictors using the available observation. No Gaussian distribution or linear regression assumption was needed.

The same geometry yields the law of total variance. Writing \(X-\mathbb E[X]=(X-M)+(M-\mathbb E[X])\) and using orthogonality,
\[
 \boxed{\operatorname{Var}(X)
 =\mathbb E[\operatorname{Var}(X\mid Y)]
  +\operatorname{Var}(\mathbb E[X\mid Y]),}
\]
where \(Y\) denotes the partition observation. Total variability splits into the uncertainty left inside each cell and the variability explained by knowing the cell.

\section{Jensen turns changes of viewpoint into inequalities}

A convex function \(\varphi\) satisfies, for nonnegative weights summing to one,
\[
 \varphi\left(\sum_i\lambda_ix_i\right)
 \le\sum_i\lambda_i\varphi(x_i).
\]
Applying this separately within each information cell gives the conditional Jensen inequality:
\[
 \boxed{\varphi(\mathbb E[X\mid Y])
        \le\mathbb E[\varphi(X)\mid Y].}
\]
This follows directly from convexity; the inequality requires the relevant expectations to exist.

The familiar Cauchy inequality
\[
 (a+b+c)^2\le3(a^2+b^2+c^2)
\]
is Jensen for \(\varphi(x)=x^2\) and a uniform random variable whose three possible values are \(a,b,c\). Similarly, concavity of \(\log\) gives weighted AM--GM:
\[
 \sum_i\lambda_i\log x_i
 \le\log\left(\sum_i\lambda_ix_i\right)
 \quad\Longleftrightarrow\quad
 \boxed{\prod_i x_i^{\lambda_i}\le\sum_i\lambda_ix_i}
\]
for \(x_i>0\). Choosing the weights and the random variable is therefore a systematic way of finding the right normalization in many inequality problems.

\section{Merging outcomes can only erase relative entropy}

Let \(p=(p_i)\) and \(q=(q_i)\) be strictly positive probability vectors on a finite set. Their relative entropy is
\[
 D(p\|q)=\sum_i p_i\log\frac{p_i}{q_i}.
\]
Define the likelihood ratio \(R(i)=p_i/q_i\) and take expectation under \(q\). Then
\[
 D(p\|q)=\mathbb E_q[R\log R].
\]
The function \(\varphi(r)=r\log r\) is strictly convex on \(r>0\), because \(\varphi''(r)=1/r\). By ordinary Jensen,
\[
 D(p\|q)\ge\varphi(\mathbb E_qR)=\varphi(1)=0,
\]
with equality precisely when \(p=q\).

Now suppose an observation merges individual outcomes into cells \(C_j\). Let
\[
 \overline p_j=\sum_{i\in C_j}p_i,\qquad
 \overline q_j=\sum_{i\in C_j}q_i.
\]
The conditional mean of \(R\) within the \(j\)-th cell, under weights \(q_i/\overline q_j\), is
\[
 \mathbb E_q[R\mid C_j]
 =\frac{\overline p_j}{\overline q_j}.
\]
Apply conditional Jensen and average over all cells:
\[
\begin{aligned}
 D(p\|q)
 &=\mathbb E_q[\varphi(R)]\\
 &\ge\sum_j\overline q_j
        \varphi\left(\frac{\overline p_j}{\overline q_j}\right)\\
 &=\sum_j\overline p_j
        \log\frac{\overline p_j}{\overline q_j}.
\end{aligned}
\]
Therefore
\[
 \boxed{D(p\|q)\ge D(\overline p\|\overline q).}
\]
This is the data-processing inequality for a deterministic coarse-graining. Distinguishing fewer outcomes cannot make two distributions more distinguishable in relative entropy. Equality holds when the ratio \(p_i/q_i\) is constant on each cell.

The projection interpretation and data processing describe the same loss of information from two angles. For squared prediction error, conditioning keeps the component visible to the observation and discards an orthogonal residual. For relative entropy, conditioning averages the likelihood ratio, and convexity shows that its distinguishing power cannot increase. Both calculations begin with reweighting on the information cells and end with a precise inequality.
